\documentclass[12pt]{article}

\usepackage{setspace}
\onehalfspace

%\usepackage[utf8]{inputenc}
%\usepackage{t1enc}
%\usepackage[T1]{fontenc}

\usepackage{amsmath}

\usepackage{moodle}

\DeclareMathOperator{\Exp}{Exp}
\DeclareMathOperator{\Norm}{{\cal N}}
\DeclareMathOperator{\Binom}{Binom}
\DeclareMathOperator{\E}{E}
\DeclareMathOperator{\D}{D}
\renewcommand{\P}{\operatorname{P}}
%\DeclareMathOperator{\cov}{cov}
\DeclareMathOperator{\corr}{corr}

\DeclareMathOperator{\F}{F} 
\renewcommand{\d}{\operatorname{d\! }} %integral \! a nyero


\DeclareMathOperator{\R}{\mathbb{R}}
%\DeclareMathOperator{\Norm}{{\cal N}}
\DeclareMathOperator{\Unif}{{\cal U}}
\DeclareMathOperator{\Poi}{{Poisson}}


\DeclareMathOperator{\cov}{cov}





\begin{document}

\begin{quiz}{gyakorló}

\begin{multi}{felt3}
$A$-nak a $B$-re vonatkozó feltételes valsége $\P(A|B)=\frac{\P(A\cdot B)}{\P(B)}$.
\item igaz
\item* hamis
\end{multi}
\begin{multi}{impdvv5}
Azt mondjuk, hogy $\xi\sim \Binom(n,p)$, ha 
$$
\P(\xi=k)=\binom{n}{k}p^k(1-p)^{n-k}\ \ \ k=0\ldots n
$$
valamilyen $n>0$ és $p\in[0,1]$ esetén.
\item* igaz
\item hamis
\end{multi}
\begin{multi}{esem2}
Az $\overline{A+B}$ és $\overline{A} \cdot \overline{B}$ események kizárják egymást.
\item igaz
\item* hamis
\end{multi}
\begin{multi}{impfvv12}
Ha $\xi\sim \Unif(a,b)$ akkor a sűrűségfüggvénye 
$$
f_{\xi}(x)=
\begin{cases}
\frac{x-a}{b-a} & \ \ a<x<b \\
0 & \text{máskor} 
\end{cases}
$$
\item igaz
\item* hamis
\end{multi}
\begin{multi}{flenvv4}
Ha $\xi$ és $\eta$ diszkrét valségi változóknak létezik a szórása és nemnulla, akkor a 
korrelációjuk
$$
\corr(\xi,\eta)=\frac{\cov(\xi,\eta)}{\D(\xi)\D(\eta)}
$$
\item* igaz
\item hamis
\end{multi}
\begin{multi}{eofv11}
Bármely $\xi$-re és $a$-ra $\P(\xi= a)=0$, feltéve hogy $F_{\xi}$ folytonos.
\item* igaz
\item hamis
\end{multi}
\begin{multi}{sfv3}
Azt mondjuk, hogy egy $f:\R\to \R$ sűrűségfüggvény 
$$
\int_{-\infty}^{+\infty}f(x)\d x\ge 0
$$
\item igaz
\item* hamis
\end{multi}
\begin{multi}{val4}
A valószínűség monoton, azaz $\P(A)\le \P(B)$, ha $A\subseteq B$.
\item* igaz
\item hamis
\end{multi}
\begin{multi}{eofv3}
Egy $F:\R\to [0,1]$ pontosan akkor eloszlásfüggvény, ha monoton nemcsökkenő, 
$\lim_{x\to+\infty}F(x)=1$ és $\lim_{x\to-\infty}F(x)=0$
\item igaz
\item* hamis
\end{multi}
\begin{multi}{impfvv3}
Egy $\xi\sim \Unif(a,b)$ szórásnégyzete: $\frac{(b-a)^2}{12}$
\item* igaz
\item hamis
\end{multi}
\begin{multi}{dvv2}
Egy nemnegatív tagú $p_1,\ldots, p_n$ vektor pont akkor valségeloszlás, ha $\sum_{k=1}^n p_k=1$.
\item* igaz
\item hamis
\end{multi}
\begin{multi}{sfv7}
Ha $\xi$-nek $f$ asűrűségfüggvénye, akkor bármely $a<b$-re:
$$
\P(a<\xi<b)=f(b)-f(a)
$$
\item igaz
\item* hamis
\end{multi}
\begin{multi}{dvv16}
$\xi$ és $\eta$ véges értékkészletű valségi változókra:
$$
\D^2(\xi+\eta)=\D^2(\xi)+\D^2(\eta)
$$
\item igaz
\item* hamis
\end{multi}
\begin{multi}{impdvv8}
Azt mondjuk, hogy $\xi\sim \Poi(\lambda)$, ha 
$$
\P(\xi=k)=e^{\lambda}\frac{\lambda^k}{k}\ \ \ (k=0,1,2\ldots )
$$ 
valamilyen $\lambda>0$ esetén.
\item igaz
\item* hamis
\end{multi}
\begin{multi}{felt10}
A teljes függetlenségből következik a páronkénti.
\item* igaz
\item hamis
\end{multi}
\begin{multi}{dvv5}
Egy $\xi$ valségi változó az $x_1,\ldots x_n$ értékeket veheti fel. Ekkor:
$$
\min_k(x_k)\le \E(\xi) \le \max_k(x_k)
$$
\item* igaz
\item hamis
\end{multi}
\begin{multi}{eofv7}
Bármely $\xi$-re és $a<b$-re $\P(a\le \xi <b)=F_{\xi}(b)-F_{\xi}(a)$.
\item* igaz
\item hamis
\end{multi}
\begin{multi}{val5}
Bármely esemény valsége kiszámolható a $\frac{\text{kedvező}}{\text{összes}}$ képlettel.
\item igaz
\item* hamis
\end{multi}
\begin{multi}{flenvv2}
Azt mondjuk, hogy $\xi$ és $\eta$ valségi változók függetlenek, ha
$$
\E(\xi \eta)=\E(\xi)\E(\eta)
$$
\item igaz
\item* hamis
\end{multi}
\begin{multi}{esem10}
Ha $\P(A)=0$ akkor $A$ sosem következik be.
\item igaz
\item* hamis
\end{multi}
\begin{multi}{sfv6}
Ha van $\xi$-nek sűrűségfüggvénye, akkor bármely $a<b$-re:
$$
\P(a<\xi<b)=\int_{a}^{b}f(x)\d x
$$
\item* igaz
\item hamis
\end{multi}
\begin{multi}{esem3}
Az $\overline{A\cdot B}$ maga után vonja az $\overline{A} + \overline{B}$ eseményt.
\item* igaz
\item hamis
\end{multi}
\begin{multi}{felt4}
$A$-nak a $B$-re vonatkozó feltételes valségét $\P(A|B)=\frac{\P(A\cdot B)}{\P(B)}$ módon értelmezzük, feltéve hogy $\P(B)>0$.
\item* igaz
\item hamis
\end{multi}
\begin{multi}{impfvv6}
Ha $\xi\sim \Unif(0,1)$ akkor  minden $a,b$-re $b\xi+a\sim \Unif(a,b)$.
\item igaz
\item* hamis
\end{multi}
\begin{multi}{dvv17}
$\xi$ és $\eta$ $\xi$ és $\eta$ véges értékkészletű valségi változókra:
$$
\D^2(\xi+\eta)=\D^2(\xi)+\D^2(\eta)+2\cov(\xi,\eta)
$$
\item* igaz
\item hamis
\end{multi}
\end{quiz}

\end{document}

